In this paper, we will approach a problem of two of the authors and Males \cite{FMR} on {\it hook polynomials} which arose in their study of the connection between zeros of partition polynomials and crank statistics. In order to explain the problem, we require hook numbers of partitions. For each cell of the Young diagram of a partition $\lambda$, the {\it hook number} of this cell is the number of boxes that form the inverted $L$-shape whose corner resides at the cell in question. For instance, the hook numbers of the partition $\lambda = \lp 5, 2, 1 \rp$ are given as
\begin{figure}[ht] ${\hspace{0.4in} \begin{ytableau} 7&5&3&2&1 \\ 3&1 \\ 1 \end{ytableau}}$
\end{figure}

 We let $\mathcal H(\lambda)$ be the multiset of all hook numbers of $\lambda$ and the $t$-hook multiset as
\begin{align*}
 \mathcal H_t(\lambda) := \{ h \in \mathcal H(\lambda) : t \mid h \}.
\end{align*}
In order to keep track of the number of $t$-hooks of partitions, we use a formula for the generating function due to Han, which derives from his extensions of the Nekrasov--Okounkov hook lengths formula \cite{Han,NekOk}. To state this theorem and those that follow, we define for each $t \geq 1$ the generating functions
\begin{align*}
 H_t\lp w;q \rp := \sum_{n \geq 0} P_t(w;n) q^n := \sum_{\lambda} w^{\#\mathcal H_t(\lambda)} q^{|\lambda|} .
\end{align*}

The following representation for $H_t(w;q)$ is proven in \cite[Corollary 5.1]{Han}.

\begin{Theorem}
 We have for $t \geq 2$ that
 \begin{align*}
 H_t(w;q) = \prod_{n \geq 1} \frac{1}{(1-(wq^t)^n)^t}\prod_{n \geq 1} \frac{\big(1-q^{tn}\big)^t}{1-q^n}.
 \end{align*}
\end{Theorem}
